% markpage — export LaTeX (cible : xelatex ou lualatex).
% La compilation avec pdflatex échouera : fontspec + UTF-8 natif
% dans les blocs de code requièrent xelatex / lualatex.
\documentclass[11pt,a4paper]{article}
\usepackage{fontspec}
\IfFontExistsTF{DejaVu Serif}{\setmainfont{DejaVu Serif}}{}
\IfFontExistsTF{DejaVu Sans}{\setsansfont{DejaVu Sans}}{}
\IfFontExistsTF{DejaVu Sans Mono}{\setmonofont{DejaVu Sans Mono}}{}

\IfFileExists{french.ldf}{\usepackage[french]{babel}}{}
\usepackage{amsmath,amssymb,amsthm}
\usepackage{stmaryrd}
\usepackage{graphicx}
\usepackage[export]{adjustbox}
\usepackage{hyperref}
\usepackage{xcolor}
\usepackage{listings}
\usepackage{enumitem}
\usepackage{booktabs}
\usepackage{tabularx}
\usepackage[normalem]{ulem}
\usepackage[breakable,skins]{tcolorbox}
\usepackage{newunicodechar}

% DejaVu Serif (our main text font) doesn't ship the Unicode
% math-bracket / logic glyphs that show up in our prose (e.g.
% explanations of the editor's ligatures). Redirect each known
% glyph to its LaTeX command via \ensuremath, so it renders
% from the math font even when typed in normal text.
\newunicodechar{⟦}{\ensuremath{\llbracket}}
\newunicodechar{⟧}{\ensuremath{\rrbracket}}
\newunicodechar{⟨}{\ensuremath{\langle}}
\newunicodechar{⟩}{\ensuremath{\rangle}}
\newunicodechar{⊢}{\ensuremath{\vdash}}
\newunicodechar{⊣}{\ensuremath{\dashv}}
\newunicodechar{⊨}{\ensuremath{\models}}
\newunicodechar{⊥}{\ensuremath{\bot}}
\newunicodechar{⊤}{\ensuremath{\top}}
\newunicodechar{∀}{\ensuremath{\forall}}
\newunicodechar{∃}{\ensuremath{\exists}}
\newunicodechar{∈}{\ensuremath{\in}}
\newunicodechar{∉}{\ensuremath{\notin}}
\newunicodechar{∅}{\ensuremath{\emptyset}}
\newunicodechar{⊂}{\ensuremath{\subset}}
\newunicodechar{⊆}{\ensuremath{\subseteq}}
\newunicodechar{⊃}{\ensuremath{\supset}}
\newunicodechar{⊇}{\ensuremath{\supseteq}}
\newunicodechar{∪}{\ensuremath{\cup}}
\newunicodechar{∩}{\ensuremath{\cap}}
\newunicodechar{∧}{\ensuremath{\land}}
\newunicodechar{∨}{\ensuremath{\lor}}
\newunicodechar{¬}{\ensuremath{\neg}}
\newunicodechar{★}{\ensuremath{\bigstar}}
\newunicodechar{✓}{\ensuremath{\checkmark}}
\newunicodechar{✗}{\ensuremath{\times}}
\newunicodechar{♥}{\ensuremath{\heartsuit}}
\newunicodechar{♦}{\ensuremath{\diamondsuit}}
\newunicodechar{♠}{\ensuremath{\spadesuit}}
\newunicodechar{♣}{\ensuremath{\clubsuit}}
\newunicodechar{⚠}{\textbf{[!]}}

\hypersetup{colorlinks=true, linkcolor=blue!50!black, urlcolor=blue!50!black}
\lstset{
  basicstyle=\ttfamily\small,
  breaklines=true,
  frame=single,
  framesep=4pt,
}

% Theorem-like environments matching the ::: admonitions (§17 of
% the markpage SPEC). The shared counter `theorem` keeps lemma /
% proposition / corollary numbered in the same sequence as the
% theorems themselves, which is the usual mathematical convention.
\newtheorem{theorem}{Théorème}
\newtheorem{lemma}[theorem]{Lemme}
\newtheorem{proposition}[theorem]{Proposition}
\newtheorem{corollary}[theorem]{Corollaire}
\theoremstyle{definition}
\newtheorem{definition}{Définition}
\newtheorem{example}{Exemple}
\newtheorem{remark}{Remarque}

\title{A complete letter — logo, sender, recipient, signature}
\author{Test Author \\ Test Org}
\date{2026-01-01}

\begin{document}
\maketitle

A full courrier in one document: a \texttt{::: background} drops the logo into
the top-right corner, while the \texttt{sender} / \texttt{recipient} / \texttt{signature}
fences handle the addressing. Exercises all four together.


\begin{tcolorbox}[breakable, colback=gray!5, colframe=gray!50, title=Background]
\textit{[image — img-logo.svg]}
\end{tcolorbox}


\begin{lstlisting}
**Atelier Méridien**
14 quai de la Mégisserie
75001 Paris
+33 1 42 33 18 90 · contact@atelier-meridien.fr
\end{lstlisting}


\begin{lstlisting}
Société Lumière & Fils
À l'attention de Mme Claire Fontaine
8 boulevard Voltaire
75011 Paris
\end{lstlisting}


Paris, le 27 juin 2026


\textbf{Objet :} Proposition d'étude — réaménagement du hall d'accueil


Madame,


Nous vous remercions de la confiance que vous accordez à notre atelier
pour le réaménagement du hall d'accueil de votre siège. Après une
première visite des lieux, nous avons le plaisir de vous adresser notre
proposition d'étude préliminaire.


Notre intervention se déroulerait en trois phases : un relevé de
l'existant et un diagnostic d'usage, une esquisse présentant deux partis
d'aménagement, puis un avant-projet détaillé chiffré.


Veuillez agréer, Madame, l'expression de nos salutations distinguées.


\begin{lstlisting}
**Julien Méridien**
*Architecte D.P.L.G.*
\end{lstlisting}

\end{document}
